\documentclass[12pt]{article}
\usepackage[margin=1in]{geometry}
\usepackage{amsmath,amsthm,amssymb,scrextend}
\usepackage{array}
\setlength{\parindent}{0pt}
\begin{document}
\newcommand{\ds}{\displaystyle}

\textbf{(i)} No ending after $2n$ tosses: $\ds P=2(pq)^n$

$$\lim_{n\to\infty}2(pq)^n=0\implies P(\text{game never ends})=0$$

Let $P_H=P(A\text{ wins}\mid \text{first toss is H}),\  P_T=P(A\text{ wins}\mid \text{first toss is T})$

$$P_H=p+qpP_H\implies P_H=\frac{p}{1-pq}$$

$$P_T=pP_H=\frac{p^2}{1-pq}$$
\begin{align*}
P(A\text{ wins})
&=pP_H+qP_T\\
&=p\cdot\frac{p}{1-pq}+q\cdot\frac{p^2}{1-pq}\\
&=\frac{p^2(1+q)}{1-pq}
\end{align*}

\bigskip\hrule\bigskip

\textbf{(ii)}

From $H$:

$$HH:\ p^2\Rightarrow A\text{ wins}$$

$$T:\ q\Rightarrow P_T,\  HT:\ pq\Rightarrow P_T$$

$$P_H=p^2+(q+pq)P_T$$

From $T$:

$$TT:\ q^2\Rightarrow B\text{ wins}$$

$$H:\ p\Rightarrow P_H,\  TH:\ pq\Rightarrow P_H$$

$$P_T=(p+pq)P_H=(1-q^2)P_H$$

$$P_H=p^2+(q+pq)(1-q^2)P_H=p^2+(1-p^2)(1-q^2)P_H$$

$$P_H=\frac{p^2}{1-(1-p^2)(1-q^2)}$$

\begin{align*}
P(A\text{ wins})
&=pP_H+qP_T\\
&=pP_H+q(1-q^2)P_H\\
&=(1-q^3)P_H\\
&=\frac{p^2(1-q^3)}{1-(1-p^2)(1-q^2)}
\end{align*}

\bigskip\hrule\bigskip

\textbf{(iii)}
$$P_H=p^{a-1}+q(1+p+\cdots+p^{a-2})P_T=p^{a-1}+(1-p^{a-1})P_T$$

$$P_T=(1-q^{b-1})P_H$$

$$P_H=p^{a-1}+(1-p^{a-1})(1-q^{b-1})P_H=\frac{p^{a-1}}{1-(1-p^{a-1})(1-q^{b-1})}$$


\begin{align*}
P(A\text{ wins})
&=pP_H+qP_T\\
&=pP_H+q(1-q^{b-1})P_H\\
&=(1-q^b)P_H\\
&=\frac{p^{a-1}(1-q^b)}{1-(1-p^{a-1})(1-q^{b-1})}
\end{align*}

When $a=b=2$,

\begin{align*}
P(A\text{ wins})
&=\frac{p(1-q^2)}{1-(1-p)(1-q)}\\
&=\frac{p^2(1+q)}{1-pq}
\end{align*}


\bigskip
\hrule
\begin{center}
    This solution was submitted by \textbf{Jon}\\[4pt]
    Published and refined by \textbf{Jon}\\[8pt]
    \textit{Last updated: July 20, 2026}
\end{center}

\end{document}